% =============================================================================
% Lecture 7: Closed-Loop Evaluation
% Physical AI: Machine Learning Systems That Sense and Act (Volume IV), Chapter 7
% =============================================================================
% Thesis: a physical success rate is conditional evidence about one machine,
% environment, and protocol, and each added nine of demonstrated reliability
% costs about ten times the trials, so evaluation states its bound and names
% the conditions it leaves to the runtime.
% Running case: the warehouse mobile manipulator's cage-door campaign,
% latch-bc-01, twenty clean runs at the 0.03 m/s guarded approach.
% Hook and synthesis: NHTSA PE 16-007 (2017), the Autosteer crash-rate
% denominator (War Story 7.1).
%
% Build: make deck V=4 D=07_evaluation   (figures: images/make_figs.py)
% Numbers: the rendered chapter; derived values from images/calcs.py.
% =============================================================================

\documentclass[aspectratio=169, 12pt]{beamer}
\usepackage{../../assets/beamerthememlsys}
\mlsyssetup{
  volume={Volume IV},
  chapter={Chapter 7},
  chaptertitle={Closed-Loop Evaluation},
  layer=world,
  logo={../../assets/img/logo-mlsysbook.png},
  instlogo={../../assets/img/logo-harvard.png},
}
\usepackage[T1]{fontenc}
\usepackage[scaled=0.9]{helvet}
\usepackage{courier}
\renewcommand{\familydefault}{\sfdefault}
\usepackage{booktabs,amsmath}
\graphicspath{{images/}}
\setcounter{mlsystotalsections}{5}

\title{Closed-Loop Evaluation}
\author{Vijay Janapa Reddi}
\institute{Harvard University}
\date{}

% A large answer choice for predict and concept-question slides
\newcommand{\choice}[2]{\textbf{\textcolor{mlsysaccent}{#1}}\enspace #2}
% A goals or takeaways line: a text-free thumbnail beside its sentence
\newcommand{\thumbline}[2]{%
  \begin{minipage}[c]{2.6cm}\centering\includegraphics[width=2.5cm,height=0.95cm,keepaspectratio]{#1}\end{minipage}%
  \hspace{0.3cm}\begin{minipage}[c]{0.74\textwidth}\raggedright\small #2\end{minipage}}
% Book notation used on the slides (books/shared/_includes/_macros.tex)
\newcommand{\asig}{\alpha_{\text{sig}}}
\newcommand{\plower}{p_{\text{lower}}}
\newcommand{\ptarget}{p_{\text{target}}}

\begin{document}

% =============================================================================
% TITLE
% =============================================================================
\mlsysnote{[0.5 min] [CORE] Read the thesis aloud once. Lecture 6 ended with a checkpoint and a policy manifest; today asks what physical trials can establish about it. The notes total about 71 minutes; the core frames take about 58. FLEX: cutting every frame tagged as optional leaves a complete lecture of about 58 minutes. Statistics carry the middle of this lecture; robotics and control students who have not seen confidence bounds lose the thread at the clean-run bound, so take that slide slowly.}
\mlsystitle{Closed-Loop Evaluation}{A success rate is conditional evidence about one machine, environment, and protocol, and each added nine of reliability costs about ten times the trials.}{cover_evaluation_blueprint.png}

% =============================================================================
% SECTION: DENOMINATORS
% =============================================================================
\mlsysnote{[0.5 min] [CORE] Open with a published safety number and ask what it divided by.}
\mlsysdivider{What does a success rate actually measure?}{Closed-Loop Evaluation, \S7.1, \S7.2, \S7.10}
\section{Denominators}

\begin{frame}{In 2017, NHTSA reported a 40 percent drop in crash rate}
  \begin{columns}[c]
    \begin{column}{0.50\textwidth}
      \centering\safeimg[width=\textwidth,height=5.0cm,keepaspectratio]{nhtsa_fig11.png}
    \end{column}
    \begin{column}{0.46\textwidth}
      {\normalsize\itshape ``The data show that the Tesla vehicles crash rate dropped by almost 40 percent after Autosteer installation.''}\\[0.4cm]
      {\small Airbag-deployment crashes per million miles, Model S and Model X, before and after Autosteer was installed.}
    \end{column}
  \end{columns}
  \mlsyscredit{NHTSA ODI, PE 16-007 closing resume (January 2017), p. 10 and Figure 11, static.nhtsa.gov/odi/inv/2016/INCLA-PE16007-7876.PDF. Closed-Loop Evaluation, War Story 7.1.}
  \note{[2.5 min] LINK: Lecture 1 named the causal boundary; a driver-assistance system crosses it every second it steers. NARRATE: in January 2017 the National Highway Traffic Safety Administration closed its investigation of a fatal Florida crash with a Model S in Autopilot mode (Lecture 8 opens on that crash itself; today uses only the report's statistic). Left: the report's own Figure 11. Airbag-deployment crashes per million miles fell from 1.3 before Autosteer was installed to 0.8 after, a 38.5 percent drop that the report rounds to ``almost 40 percent'' (right, the report's own sentence). The number was quoted widely as evidence that the lane-keeping system made the cars safer. ENGAGE: before the next slide, ask the class to write down what a crash rate needs in its denominator for that conclusion to hold. Expected answer: miles driven with Autosteer engaged, for the same vehicles, over the same kinds of road. WARN: the bars are real and correctly drawn; the problem is not arithmetic. FLEX: [CORE] the hook; keep it.}
\end{frame}

\begin{frame}{Which miles did the 0.8 crashes per million divide by?}
  \mlsysactivity{Predict}{2 min}
  \centering
  {\normalsize\choice{A}{Miles driven with Autosteer engaged}\qquad\choice{B}{All miles after installation}}\\[0.5cm]
  \uncover<2->{%
    {\normalsize Footnote 22: the rates cover \textbf{all miles} before and after installation,\\not only miles driven with Autopilot in use.}\\[0.4cm]
    {\normalsize Cohort with complete mileage: 0.76 $\rightarrow$ 1.21 crashes per million miles, \textbf{+59\%}.}}
  \mlsyscredit{Footnote 22: NHTSA PE 16-007, p. 10. Replication: Quality Control Systems Corp. 2019, via Closed-Loop Evaluation, War Story 7.1.}
  \note{[2 min] ENGAGE: hands up for A, then for B. Most of the room picks A. NARRATE, step 2: the report's own footnote 22 says the rates cover all miles driven before and after Autopilot installation, not only miles with Autopilot in use. The denominator never measured exposure to Autosteer. Two years later, after a Freedom of Information Act lawsuit released the per-vehicle mileage, Quality Control Systems Corporation found that exact exposure before and after installation could be computed for only 13 percent of vehicles. In the one cohort with complete exposure, the rate rose 59 percent, from 0.76 to 1.21 per million miles. WARN: the replication does not measure Autosteer's effect either; it shows the published denominator could not support the reported drop. Say the line: ``Today is about denominators: what was sampled, what was counted, and what the count can bound.'' FLEX: [CORE].}
\end{frame}

\begin{frame}{Recall: what does a policy manifest declare?}
  \mlsysactivity{Retrieval}{1.5 min}
  \centering
  {\Large\bfseries Close your notes.}\\[0.2cm]
  {\normalsize Write what latch-bc-01's manifest declares as its supported region,\\and what it records about the simulator.}\\[0.4cm]
  \uncover<2->{%
    \thumbline{thumb_coverage_step2.pdf}{One cell of six, the nominal strike plate at the 0.03 m/s guarded approach. An omission log lists what training never saw.}}
  \mlsyscredit{Closed-Loop Evaluation, \S7.1 and \S7.8; the policy manifest is defined in Policy Training.}
  \note{[1.5 min] LINK: Lecture 6 built the policy manifest. Give 45 seconds. NARRATE, step 2: the checkpoint latch-bc-01 came from offline imitation of hardware demonstrations. Its manifest declares one supported condition, the nominal strike plate at the guarded 0.03 m/s approach, one cell of a six-cell scenario table (three plate positions by two approach speeds). The manifest also carries an omission log of what the training never modeled. The thumbnail is today's coverage grid; it returns in segment 1. ENGAGE: ``If the manifest supports one cell, what can trials in that cell add?'' Expected: a bound on success in that cell, and nothing about the other five. FLEX: [CORE] a student who missed Lecture 6 needs only this slide.}
\end{frame}

\begin{frame}{Why can a robot pass every trial and still not be ready?}
  \thumbline{thumb_coverage_step2.pdf}{State what a physical success rate measures, and keep every intervention out of the policy's credit.}\\[0.22cm]
  \thumbline{thumb_regimes_step4.pdf}{Choose the evidence regime that can expose a given failure: replay, simulation, shadow, or a constrained physical test.}\\[0.22cm]
  \thumbline{thumb_nines_step4.pdf}{Compute the bound a clean record supports and the trials a target needs.}\\[0.22cm]
  \thumbline{thumb_record_step2.pdf}{Price the exposure wall, and write the evaluation record that hands untested conditions to the runtime.}\\[0.15cm]
  {\footnotesize\textcolor{midgray}{Thread: the second law, endogenous data. Record: the evaluation record.}}
  \note{[1.5 min] NARRATE: the title is the chapter's driving question. Four goals, one per segment. LINK: Lecture 5 stated the second law, endogenous data, the thread on the bottom line: a robot's actions choose the states it will meet. Today it acts on testing, and the lecture closes Part II by returning to it. A test campaign samples the trajectories that the policy and the permission path produce, so no recorded dataset can stand in for them. The lecture adds one record to the record chain: the evaluation record, the link after the policy manifest. Each thumbnail is the shape of the figure that teaches its goal. WARN: the second law is a design principle the book elevates to a law; only the first law, irreversibility, is physical. FLEX: [CORE].}
\end{frame}

\begin{frame}{The running case is a mobile manipulator opening a door}
  \begin{columns}[c]
    \begin{column}{0.44\textwidth}
      \centering\safeimg[width=\textwidth,height=5.0cm,keepaspectratio]{erflex.jpg}\\[0.05cm]
      {\footnotesize MiR250 base carrying a UR16e arm}
    \end{column}
    \begin{column}{0.52\textwidth}
      \small
      \begin{mlsyscard}{mlsysaccent}
        \textbf{Campaign for latch-bc-01}\\[0.1cm]
        Base stops 2 m out; arm moves at 0.03 m/s\\
        Abort if contact force passes the 15 N tripwire\\
        Latch structural limit: 100 N\\
        One run 25 s; reset 15 s\\[0.1cm]
        \textbf{Result: 20 runs, 20 clean.}
      \end{mlsyscard}
    \end{column}
  \end{columns}
  \mlsyscredit{Photo: ER-FLEX by Enabled Robotics; Auledas, CC BY 4.0, Wikimedia Commons. Campaign: Closed-Loop Evaluation, \S7.1 and \S7.7.2 (illustrative values).}
  \note{[1.5 min] NARRATE: the book's running machine is a reference platform, a wheeled base with an articulated arm. The photo is a real machine of that class, a commercial mobile robot carrying a collaborative arm; the book's machine is a reference design of the same class. Walk the card. In each trial the base drives in from a 2 m standoff and stops; the arm approaches the spring-latched cage door at a guarded 0.03 m/s, unlatches it, and swings it open without crossing a 15 N contact tripwire or drawing a human intervention. A run takes 25 s and an automated reset 15 s. The team ran twenty trials and all twenty were clean. ENGAGE: ``On a vision benchmark, 100 percent on a held-out set would ship. What is missing here?'' Collect two answers; the lecture answers them. WARN: these campaign values are illustrative in the book; the statistics that follow do not depend on them. FLEX: [CORE].}
\end{frame}

\begin{frame}{A success rate belongs to a population of trial conditions}
  \centering
  {\small\textcolor{midgray}{The closed-loop estimand}}\\[0.15cm]
  {\Large $p(\mathcal{D}) \;=\; \mathbb{E}_{c\,\sim\,\mathcal{D}}\bigl[\,\phi(\text{Traj}(c))\,\bigr]$}\\[0.3cm]
  {\large $c = (\,
    \uncover<2->{\textcolor{BlueLine}{\mathbf{x}_0}},\;
    \uncover<2->{\textcolor{GreenLine}{e}},\;
    \uncover<3->{\textcolor{OrangeLine}{m}},\;
    \uncover<4->{\textcolor{VioletLine}{\pi_\theta}},\;
    \uncover<4->{\textcolor{midgray}{\kappa}}\,)$}\\[0.25cm]
  \small
  \begin{tabular}{@{}P{0.175\textwidth}P{0.175\textwidth}P{0.175\textwidth}P{0.175\textwidth}P{0.175\textwidth}@{}}
    \uncover<2->{\textcolor{BlueLine}{\textbf{start state}}} &
    \uncover<2->{\textcolor{GreenLine}{\textbf{environment}}} &
    \uncover<3->{\textcolor{OrangeLine}{\textbf{this unit}}} &
    \uncover<4->{\textcolor{VioletLine}{\textbf{frozen policy}}} &
    \uncover<4->{\textcolor{midgray}{\textbf{protocol}}} \\
    \uncover<2->{pose, speed} &
    \uncover<2->{friction, light, heat} &
    \uncover<3->{gear play, battery, wear} &
    \uncover<4->{weights, precision} &
    \uncover<4->{deadline, goal, predicate} \\
  \end{tabular}\\[0.1cm]
  {\normalsize\uncover<5->{Change friction or battery voltage and $p$ changes with no weight touched.}}
  \mlsyscredit{Closed-Loop Evaluation, \S7.2, equation 1 and Definition 7.1 (closed-loop evaluation estimand).}
  \note{[3 min] NARRATE: the chapter's estimand. $p$ is the expected value of a frozen pass-fail predicate $\phi$ applied to the trajectory a trial condition $c$ produces, averaged over a declared distribution $\mathcal{D}$ of conditions. Step 2: the start state and the environment: floor friction, light, air temperature. Step 3: the unit itself. Nominally identical robots differ in transmission backlash, camera mounting, torque constants, and battery charge (rig-to-rig divergence, Lecture 5). Step 4: the frozen policy, with its numerical precision and the permission path's parameters, and the protocol: deadline, goal region, and the predicate. Step 5: the punchline. ENGAGE: ``A worn tire. Did the model change?'' Expected: no weight changed, but the plant transition changed, so $p$ changed. WARN: machine learning students treat held-out accuracy as a property of the weights; here it is a property of the whole closed loop and the population sampled. Any policy can score 100 percent if tested only in benign states. FLEX: [CORE] concept trace: closed-loop estimand.}
\end{frame}

\begin{frame}{Twenty clean runs sampled one cell of six}
  \begin{columns}[c]
    \begin{column}{0.54\textwidth}
      \centering\mlsyssteps[width=\textwidth,height=5.3cm,keepaspectratio]{coverage_step}{2}
    \end{column}
    \begin{column}{0.42\textwidth}
      \centering\safeimg[width=\textwidth,height=4.4cm,keepaspectratio]{tri_settings.png}\\[0.05cm]
      {\footnotesize TRI's planned start scenes: nominal, object shift, station shift.}
    \end{column}
  \end{columns}
  \mlsyscredit{Left: Closed-Loop Evaluation, \S7.8 and \S7.10.1 (Pathology 1). Right: TRI LBM Team 2025, Figure 10 (evaluation settings), arXiv 2507.05331.}
  \note{[2 min] NARRATE, step 1: the scenario table crosses three strike-plate positions with two approach speeds. The retained demonstrations covered one cell, the nominal plate at the guarded speed; the operational design domain (ODD) of Lecture 5 is declared from it. Step 2: all twenty runs came from that cell, with the base stopped, the tote rack empty, on Unit 01. Each cell, or stratum, is its own population with its own success probability; a claim about the whole table needs weights declared before the first run. Right: a real evaluation design. The Toyota Research Institute team fixed three settings before evaluating its policies: nominal, object-centric shift (new objects), and station shift (a different station); each row shows real camera frames of planned start scenes. ENGAGE: ``The report says the robot opens the site's cage doors. Which cells support that?'' Expected: one. WARN: the chapter calls this environment selection, the test bench narrower than the claim. FLEX: [CORE].}
\end{frame}

\begin{frame}{Completion, autonomy, and burden are separate outcome channels}
  \centering
  \small
  \renewcommand{\arraystretch}{1.35}
  \begin{tabular}{@{}lccc@{}}
    \toprule
    \textbf{What happened in the door run} & \textbf{Task done $Y$} & \textbf{Unassisted $A$} & \textbf{Interventions} \\
    \midrule
    \uncover<1->{Policy opens the door alone} & \uncover<1->{\textcolor{GreenLine}{\textbf{1}}} & \uncover<1->{\textcolor{GreenLine}{\textbf{1}}} & \uncover<1->{0} \\
    \uncover<2->{Permission path clips a high force; door opens} & \uncover<2->{\textcolor{GreenLine}{\textbf{1}}} & \uncover<2->{\textcolor{RedLine}{\textbf{0}}} & \uncover<2->{1 clip} \\
    \uncover<3->{Permission path halts the arm at a joint limit} & \uncover<3->{\textcolor{RedLine}{\textbf{0}}} & \uncover<3->{\textcolor{RedLine}{\textbf{0}}} & \uncover<3->{1 stop} \\
    \uncover<4->{Operator clears a jammed pin; robot finishes} & \uncover<4->{\textcolor{GreenLine}{\textbf{1}}} & \uncover<4->{\textcolor{RedLine}{\textbf{0}}} & \uncover<4->{1 human} \\
    \bottomrule
  \end{tabular}\\[0.35cm]
  {\small\uncover<5->{ACR (autonomous completion rate) counts only $A = 1$; MTBI (mean time between interventions) prices the burden.\\Intervened trials are censored, so their unassisted outcome stays unknown.}}
  \mlsyscredit{Closed-Loop Evaluation, \S7.2.2 (four outcomes; ACR, MTBI, intervention density, reset latency).}
  \note{[3 min] NARRATE: four door runs, one row each. Step 1: the policy acts alone; the run counts on every channel. Step 2: the policy proposes too much force and the permission path clips it; the door opens. The whole machine completed the task ($Y = 1$), and the proposer did not do it unassisted ($A = 0$). Step 3: the permission path halts the arm; hardware protected, task failed. Step 4: an operator clears a jammed pin; the task completes on human labor. Step 5: three channels: task completion $Y$, the autonomous completion rate (ACR, runs with $Y=1$ and zero interventions), and supervisory burden, measured as mean time between interventions (MTBI) or interventions per kilometer. LINK: Lecture 1's proposal boundary; only the first row credits the proposer. WARN: a clip is the architecture working as designed, so it must not brand the run a system failure; it still counts against autonomy. A jammed run's 60 s manual reset is 70.6 percent of that run's cell time (derived, 60 of 85 s). FLEX: [CORE] concept trace: outcome channels.}
\end{frame}

\begin{frame}{Vote: an operator stops a wedging gripper. Score it?}
  \mlsysactivity{Vote}{2 min}
  \centering
  {\normalsize Run 10 of 10: the gripper wedges against the latch, and the operator takes over before anything breaks.}\\[0.35cm]
  {\normalsize\choice{A}{Discard it: an invalid trial}\quad\choice{B}{A success: nothing broke}\quad\choice{C}{A failure for autonomy}}\\[0.45cm]
  \uncover<2->{%
    \begin{mlsyscard}{mlsysaccent}
      \small\textbf{C.} It stays in the denominator, so the rate is 9 of 10, 90 percent. Dropping it reports 100 percent and an infinite MTBI.\\[0.05cm]
      Whether it would have crashed is unknown, so the run is \textbf{right-censored}.
    \end{mlsyscard}}
  \mlsyscredit{Closed-Loop Evaluation, \S7.4.1 (censored trial fallacy) and \S7.10.1 (Pathology 5, denominator filtering).}
  \note{[2 min] ENGAGE: vote by hands, then thirty seconds with a neighbor, then vote again. Expected spread: many pick A, because the run ``did not finish''. Step 2: C. The run answers two questions, and the record keeps them apart. For autonomy it is an unambiguous failure, $A = 0$. For the separate question of whether the unassisted policy would have damaged something, the intervention ended the trial first, so the outcome is unknown; survival analysis calls this right-censored (Kaplan and Meier 1958). Dropping it moves every metric the flattering way. ACR rises from 90 to 100 percent and MTBI becomes infinite, with no weight changed. Crediting the policy with the protection the operator supplied is the censored trial fallacy. LINK: Lecture 5's intervention curation keeps these runs as data for the same reason. WARN: B is the zero-damage fallacy; zero damage under supervision measures the supervisor. Land the segment: a success rate is a count over a declared population, with every intervention kept and attributed. FLEX: [CORE] concept trace: censored trials.}
\end{frame}

% =============================================================================
% SECTION: CLOSING THE LOOP
% =============================================================================
\mlsysnote{[0.5 min] [CORE] Before any physical trial, a team can score the policy against logs. Why is that not enough?}
\mlsysdivider{Why can a low offline error drive into a rack?}{Closed-Loop Evaluation, \S7.3 and \S7.4}
\section{Closing the Loop}

\begin{frame}{The lowest validation loss picked a 2.7 percent policy}
  \begin{columns}[c]
    \begin{column}{0.50\textwidth}
      \centering\safeimg[width=\textwidth,height=5.2cm,keepaspectratio]{mandlekar_bars.pdf}
    \end{column}
    \begin{column}{0.46\textwidth}
      \centering\safeimg[width=\textwidth,keepaspectratio]{mandlekar_g1.png}\\[0.1cm]
      {\footnotesize Success rate (left) climbs while validation loss (right) climbs too.}
    \end{column}
  \end{columns}
  \mlsyscredit{Mandlekar et al., CoRL 2021, arXiv 2108.03298, Appendix G and Figure G.1(a, b): BC-RNN, low-dim, 30\% validation split, 3 seeds. Closed-Loop Evaluation, \S7.3.}
  \note{[2 min] NARRATE: published data. Mandlekar and colleagues trained imitation policies on robot manipulation demonstrations and rolled every checkpoint out in closed loop. Left: on the Square task, the highest-scoring checkpoint succeeded 80.7 percent of the time; the checkpoint with the lowest validation loss succeeded 2.7 percent. On Transport, 64.0 against 0.7 percent. Right: their own curves for Square. Validation loss bottoms out early, near epoch 100 to 300, and then rises, while closed-loop success keeps climbing. The offline evaluation illusion is the assumption that lower loss on logged states means better execution. ENGAGE: ``Validation loss is computed on held-out demonstrations. Why would it disagree with rollouts?'' Expected: it scores each step from the demonstrator's state, so it never sees the states the policy's own errors lead to. LINK: compounding error from Lecture 5. WARN: offline metrics still screen out broken checkpoints and check latency. FLEX: [CORE] concept trace: offline evaluation illusion.}
\end{frame}

\begin{frame}{A 1° heading bias: how far down a 30 m aisle?}
  \mlsysactivity{Predict}{3 min}
  \centering
  {\small Base at 1.5 m/s, 15 cm clearance each side. Replay scores the bias at MSE $3.05\times10^{-4}$ rad$^2$.}\\[0.15cm]
  {\normalsize\choice{A}{Never hits}\quad\choice{B}{About 25 m}\quad\choice{C}{About 9 m}\quad\choice{D}{About 1 m}}\\[0.2cm]
  \only<1|handout:0>{\safeimg[width=0.95\textwidth,height=4.6cm,keepaspectratio]{aisle_step1.pdf}}%
  \only<2-|handout:1>{\safeimg[width=0.95\textwidth,height=4.6cm,keepaspectratio]{aisle_step2.pdf}}
  \mlsyscredit{Closed-Loop Evaluation, \S7.3.1, equation 2 and figure 1 (navigation stack; values derived in images/calcs.py).}
  \note{[3 min] NARRATE, step 1: the navigation policy carries a constant 1 degree heading bias. In replay the harness feeds the next recorded state, always on the centerline, so the error never accumulates, and the per-step score stays at $3.05\times10^{-4}$ rad$^2$ however long the aisle is. ENGAGE: students vote, then explain to a neighbor. Step 2: in closed loop the base actually moves. Under the unicycle model lateral drift grows as $v\,\sin\Delta\psi\,t$; at 1.5 m/s that spends the 15 cm side clearance in $0.15/(1.5\sin 1^\circ) = 5.7$ s, after 8.6 m, less than a third of the aisle (derived). Answer C. The point on the plot is the collision horizon; a permission path must act earlier, allowing detection and braking. WARN: the small number comes from radians squared. Physical causes such as tire wear or floor camber produce exactly this kind of bias. A yaw-rate bias would make drift grow with the square of time. FLEX: [CORE] concept trace: offline evaluation illusion.}
\end{frame}

\begin{frame}{At a stiff latch, 2 mm of error is 1,000 N}
  \centering
  \safeimg[width=0.92\textwidth,height=4.6cm,keepaspectratio]{contact_force.pdf}\\[0.15cm]
  {\normalsize Mean squared error scores 2 mm in free space and 2 mm into a latch the same.}
  \mlsyscredit{Closed-Loop Evaluation, \S7.3.2 (contact stiffness $5\times10^{5}$ N/m; illustrative). Penetrations at 15 N and 100 N derived.}
  \note{[1.5 min] NARRATE: contact makes the dynamics discontinuous. Against a stiff interface of $5\times10^{5}$ N/m, force grows as stiffness times penetration. The 15 N tripwire trips at 0.03 mm, the 100 N latch limit is reached at 0.2 mm, and a 2 mm residual would demand 1,000 N (derived; real loads depend on compliance and torque limits). In free space the same 2 mm is harmless, and mean squared error scores both identically. Cosine similarity fails the other way, because it ignores magnitude, so a proposal with 0.999 cosine and 25 percent too little force stalls against the latch's preload. ENGAGE: ``Two checkpoints: A has lower MSE concentrated at contact, B has larger errors in free space. Which opens more doors?'' Expected: B; the chapter calls this a rank inversion. FLEX: [OPTIONAL] cut if short; the aisle example carries the segment.}
\end{frame}

\begin{frame}{Each evidence regime trades one blind spot for another}
  \centering
  \mlsyssteps[width=0.98\textwidth,height=5.9cm,keepaspectratio]{regimes_step}{4}
  \mlsyscredit{Closed-Loop Evaluation, \S7.4.1, figure 2 and table 2, redrawn at slide size.}
  \note{[3 min] NARRATE: two axes. Horizontal, causal feedback: do the candidate's actions change its next input? Vertical, physical realism: do observations come from the real machine? Step 1: recorded replay. Good for latency, memory footprint, and numerics; blind to drift, because errors never feed back. Step 2: closed-loop simulation restores feedback in a model. It tests recovery across thousands of parallel runs, and its evidence applies to the model; backlash, tire slip, and real contact are absent by construction. Step 3: shadow operation runs the candidate on the real machine while a validated baseline drives. It sees real buses, heat, and sensor noise. Because the baseline keeps authority, the candidate only sees states the baseline visits, so its recovery is never tested, which the chapter calls trajectory confinement. LINK: Lecture 16 uses shadow mode as a release filter. Step 4: constrained physical testing closes the loop on hardware inside a safety envelope (red dashed); it censors everything beyond that boundary. WARN: students read the regimes as a ladder of fidelity; each step exchanges one distortion for another. FLEX: [CORE] concept trace: four evidence regimes.}
\end{frame}

\begin{frame}{A simulator earns screening authority from paired real runs}
  \centering
  \safeimg[width=0.95\textwidth,height=4.7cm,keepaspectratio]{simpler_fig1.png}\\[0.1cm]
  {\normalsize Same policies run real and simulated, about 1,500 paired episodes.}
  \mlsyscredit{Li et al., SIMPLER, CoRL 2024, Figure 1, arXiv 2405.05941. Closed-Loop Evaluation, \S7.4.2.}
  \note{[2 min] NARRATE: a real published evaluation. Top left, real-robot evaluation: expensive, slow, hard to reproduce. Bottom left, the SIMPLER team built simulated scenes matched to two real robot setups. Right: each point is one policy on one task, real success on the horizontal axis, simulated success on the vertical, across about 1,500 evaluation episodes. Points near the diagonal mean the simulator ranks policies as the real cell does. The chapter's rule is that a simulator earns the authority to screen candidates only after a paired benchmark like this, scored not only by correlation but by rank inversions (the paper's mean maximum rank violation), because a high correlation can hide one costly swap. ENGAGE: ``Policy A beats B in a validated simulator. Does A meet a 99 percent target?'' Expected: unknown; ranking says nothing about absolute failure rates. WARN: the omission log in the policy manifest lists what the simulator cannot claim. FLEX: [OPTIONAL] simulator screening validity, a secondary concept; if cut, say the rule in one sentence on the next slide.}
\end{frame}

\begin{frame}{Which regime can answer each question?}
  \mlsysactivity{Your turn}{2 min}
  \small
  \renewcommand{\arraystretch}{1.2}
  \begin{tabular}{@{}>{\raggedright\arraybackslash}m{0.60\textwidth}>{\raggedright\arraybackslash}m{0.34\textwidth}@{}}
    \toprule
    \textbf{Question about the latch policy} & \uncover<2->{\textbf{Regime}} \\
    \midrule
    Does inference hold 20 Hz under heat and bus load? & \uncover<2->{Shadow operation} \\
    Can the arm recover from a wide range of positional offsets? & \uncover<2->{Closed-loop simulation} \\
    Does peak contact force stay under 100 N at the real latch? & \uncover<2->{Constrained physical test} \\
    Does inference fit its memory and time budget at all? & \uncover<2->{Recorded replay} \\
    \bottomrule
  \end{tabular}
  \mlsyscredit{Closed-Loop Evaluation, \S7.4.1 (assigning questions to regimes) and \S7.3.3.}
  \note{[2 min] ENGAGE: pairs assign a regime to each row, ninety seconds. Step 2: the chapter's assignments. Inference rate under real thermal throttling and fieldbus contention needs the real machine while the candidate holds no authority, so shadow operation answers it. Recovery from many offsets needs closed loop and many runs without hardware wear, so simulation answers it, and its result qualifies a checkpoint for physical trials. Contact force at a real latch needs real contact physics, so it needs a constrained physical test. Fitting the budget at all needs only the computation, so replay comes first. WARN: the robot's two processors matter here. Replay on the application processor cannot see the bus to the safety microcontroller, a late chunk arriving after its chunk lease has lapsed (Lecture 4), or a clip that changes the motion. Land the segment: open-loop scores cannot certify closed-loop behavior, and every regime is blind to something. FLEX: [CORE].}
\end{frame}

% =============================================================================
% SECTION: CLEAN RECORDS
% =============================================================================
\mlsysnote{[0.5 min] [CORE] Twenty clean runs on the real latch. What do they prove?}
\mlsysdivider{How much reliability do twenty clean runs show?}{Closed-Loop Evaluation, \S7.5 to \S7.7}
\section{Clean Records}

\begin{frame}{A protocol fixes draws, predicate, and denominator first}
  \centering
  \safeimg[width=0.95\textwidth,height=2.9cm,keepaspectratio]{tri_initial_overlay.png}\\[0.05cm]
  {\footnotesize Operator view: the planned start (middle) overlaid on the live camera (left) before each run.}\\[0.3cm]
  \small
  \begin{tabular}{@{}lll@{}}
    \textbf{Independent draws} & \textbf{Frozen predicate} & \textbf{Full denominator} \\
    cool-downs, recharges, fresh units & no mid-campaign edits & every abort and takeover \\
  \end{tabular}
  \mlsyscredit{TRI LBM Team 2025, Figure 8, arXiv 2507.05331. Conditions: Closed-Loop Evaluation, \S7.5; reporting practice: Kress-Gazit et al. 2024.}
  \note{[2.5 min] NARRATE: a real tool. Before every hardware trial, TRI's operators see the planned initial condition overlaid on the live camera image and arrange the scene to match it, so the start states are drawn from a declared distribution instead of chosen by the operator. The chapter's three conditions follow. Independent draws: heat in the windings, battery sag, and wear correlate run 20 with run 1, so the protocol enforces cool-downs and recharges, and twenty runs on Unit 01 bound Unit 01 only. A frozen predicate: no widening the goal region or the deadline mid-campaign. A full denominator: communication faults, timeouts, protective stops, and takeovers all stay in the log. Fix all three, plus the target ODD and its weights, before the first run. ENGAGE: ``An operator resets the base to an easier pose after a hard run. Which condition breaks?'' Expected: independent draws from the declared distribution. WARN: Kress-Gazit and colleagues made this a reporting standard: run counts, initial conditions, success criteria, and failure analysis. FLEX: [CORE] concept trace: predeclared trial protocol.}
\end{frame}

\begin{frame}{Twenty clean runs still allow one failure in seven}
  \centering
  \mlsyssteps[width=0.92\textwidth,height=4.5cm,keepaspectratio]{inversion_step}{2}\\[0.1cm]
  {\normalsize\uncover<2->{$p^{20} = \asig$ gives $\plower = 0.05^{1/20} = $ \textbf{0.861}, so failure may reach 13.9 percent.}}
  \mlsyscredit{Closed-Loop Evaluation, \S7.6, figure 3, equation 3 and Napkin Math 7.1; Clopper and Pearson 1934. Values from images/calcs.py.}
  \note{[3 min] NARRATE: the textbook interval, $\hat{p} \pm 1.96\sqrt{\hat{p}(1-\hat{p})/n}$, has zero width at $\hat{p} = 1$; the normal approximation breaks at the boundary. Use the product rule instead. Step 1: the probability that 20 independent runs are all clean is $p^{20}$, plotted for every candidate true success rate $p$. Step 2: a policy with $p$ below 0.861 would produce this clean record less than once in twenty campaigns, so the record rules those out at 95 percent confidence. That is the exact one-sided Clopper-Pearson lower limit: 86.1 percent success, a failure rate up to 13.9 percent, about one run in seven. The orange point is a policy that fails one run in eight; it still logs twenty clean runs in 6.9 percent of campaigns (derived). ENGAGE: ``Is there a 95 percent chance this policy exceeds 86.1 percent?'' Expected: no; 95 percent describes the procedure over repeated campaigns. WARN: control and robotics students who skipped statistics stall here; say the sentence twice. FLEX: [CORE] concept trace: zero-failure bound.}
\end{frame}

\mlsysnote{[1 min] [CORE] The key law of the lecture. Expand $\ln(1-p_{\text{fail}}) \approx -p_{\text{fail}}$ and set $(1 - p_{\text{fail}})^n = 0.05$: $n\,p_{\text{fail}} \approx \ln 20 = 2.996 \approx 3$. The upper limit on the failure rate after $n$ clean trials is about $3/n$, the rule of three from clinical trials (Hanley and Lippman-Hand 1983). For 20 runs, $3/20 = 0.15$ against the exact 13.9 percent. Inverted, a target needs $n \approx 3/(1 - p_{\text{target}})$ runs, so each added nine multiplies the runs by ten. WARN: the numerator, about 3, is fixed by the confidence level; only more trials improve the bound.}
\mlsysfocus{A nine costs ten times the trials}{%
  {\Large $n \;=\; \left\lceil \dfrac{\ln \asig}{\ln \ptarget} \right\rceil \;\approx\; \dfrac{3}{1 - \ptarget}$}\\[0.6cm]
  {\large 90\% $\rightarrow$ 29 \qquad 99\% $\rightarrow$ 299 \qquad 99.9\% $\rightarrow$ 2,995}}

\begin{frame}{A 99.9 percent claim needs 2,995 runs, or 33.3 hours}
  \begin{columns}[c]
    \begin{column}{0.64\textwidth}
      \centering\mlsyssteps[width=\textwidth,height=4.3cm,keepaspectratio]{nines_step}{4}
    \end{column}
    \begin{column}{0.33\textwidth}
      \centering\safeimg[width=\textwidth,height=3.6cm,keepaspectratio]{nist-stairs.jpg}\\[0.05cm]
      {\footnotesize Every run occupies a test cell and its operators (NIST stair station).}
    \end{column}
  \end{columns}
  \mlsyscredit{Closed-Loop Evaluation, Napkin Math 7.2 and table 3 (40 s door cycle, illustrative). Photo: NIST, Public domain, Wikimedia Commons.}
  \note{[2 min] NARRATE: each bar is the clean-run count a 95 percent lower limit needs, with the cell time at the 40 s door cycle. Step 1: the campaign's twenty runs, 86.1 percent, 13.3 minutes. Step 2: 90 percent needs 29 runs. Step 3: 99 percent needs 299, 3.3 hours. Step 4: 99.9 percent needs 2,995 clean runs, 33.3 hours of continuous cell time and 66.6 operator-hours with two safety operators, about eight shifts (derived). Right: a real test cell, a bomb-squad member driving a robot through NIST's standardized stair station; every trial costs that floor space and those people. ENGAGE: ``What happens to the arm over 2,995 runs?'' Expected: wear; backlash, damping, and winding heat drift, so the campaign outlives the stationarity it assumes. A component replaced midway may start a new population. LINK: Lecture 5 budgeted demonstration wear the same way. FLEX: [CORE] concept trace: zero-failure bound.}
\end{frame}

\begin{frame}{Estimate: what does 99.5 percent cost?}
  \mlsysactivity{Estimate}{2.5 min}
  \begin{columns}[T]
    \begin{column}{0.40\textwidth}
      \begin{mlsyscard}{mlsysaccent}
        \small\textbf{Given}\\[0.05cm]
        Target $\ptarget = 0.995$\\
        Confidence 95\%, zero failures\\
        Door cycle 40 s
      \end{mlsyscard}
    \end{column}
    \begin{column}{0.56\textwidth}
      \normalsize
      How many clean runs, and how many cell hours?\\[0.3cm]
      \uncover<2->{Rule of three: $3 / 0.005 = $ \textbf{600} runs}\\[0.15cm]
      \uncover<3->{Exact: $\lceil \ln 0.05 / \ln 0.995 \rceil = $ \textbf{598} runs}\\[0.15cm]
      \uncover<4->{$598 \times 40$ s $=$ \textbf{6.6 h} of cell time}
    \end{column}
  \end{columns}
  \mlsyscredit{Closed-Loop Evaluation, Checkpoint 7.2 (trial-budget estimate). Values from images/calcs.py.}
  \note{[2.5 min] ENGAGE: individual work, ninety seconds, napkin only. Step 2: the rule of three: failure rate 0.005, so about $3/0.005 = 600$ runs. Step 3: the exact count from equation 5 is 598, because $\ln(0.995)$ is slightly larger in magnitude than 0.005. Step 4: 598 runs at 40 s each take 6.6 hours of cell time. Then ask: ``Run 412 fails. Can you restart the counter and keep going?'' Expected: no. Restarting after a failure is optional stopping; the record must apply the bound with failures instead. With one failure in twenty runs, the lower limit drops from 86.1 to 78.4 percent (equation 6, derived), even though the observed rate is 95 percent. WARN: students multiply by 3 and forget the 1 minus; write the failure rate first. FLEX: [CORE].}
\end{frame}

\begin{frame}{Stop early only on a boundary declared in advance}
  \centering
  \safeimg[width=0.95\textwidth,height=4.3cm,keepaspectratio]{sprt.pdf}\\[0.1cm]
  {\small The log-likelihood ratio scores the evidence: $+0.0953$ per clean run, $-2.303$ per failure.\\Both error rates and both boundaries are declared before run 1.}
  \mlsyscredit{Closed-Loop Evaluation, \S7.7.3, equations 7 and 8 (illustrative $p_0 = 0.10$, $p_1 = 0.01$, $\asig = 0.05$, $\beta = 0.10$); Wald 1945.}
  \note{[2 min] NARRATE: watching runs and stopping when the tally looks good inflates false acceptance; five independent 5 percent looks give a 22.6 percent chance of at least one false positive (derived). Wald's sequential probability ratio test (SPRT) fixes the stopping rule first. It compares an unacceptable failure rate, 10 percent, with an acceptable one, 1 percent. Each clean run adds 0.0953 to the log-likelihood ratio; each failure subtracts 2.303. The campaign accepts at 2.89 and rejects at $-2.25$, both set from the declared error rates. A clean record accepts after 31 runs. After 25 clean runs and one failure the score sits at 0.08 and testing continues. ENGAGE: ``A fixed test needs 29 clean runs for 90 percent. Why prefer 31?'' Expected: the fixed test rejects an acceptable 1-percent policy in 25.3 percent of campaigns; the sequential test bounds that error too. FLEX: [OPTIONAL] sequential tests; cut if short.}
\end{frame}

\begin{frame}{20 of 20 against 16 of 20 does not separate two policies}
  \begin{columns}[c]
    \begin{column}{0.54\textwidth}
      \centering\safeimg[width=\textwidth,height=4.6cm,keepaspectratio]{tri_platform.png}\\[0.05cm]
      {\footnotesize TRI ran 1,800 blind A/B rollouts.}
    \end{column}
    \begin{column}{0.42\textwidth}\raggedright
      \small
      Interleave, randomize, and blind the runs on one unit.\\[0.3cm]
      Fisher's exact test, 20/20 against 16/20: $p = $ \textbf{0.106}.\\[0.3cm]
      90 against 80 percent at 80\% power: \textbf{195} runs per arm.
    \end{column}
  \end{columns}
  \mlsyscredit{Closed-Loop Evaluation, \S7.7.4 (values derived in images/calcs.py). Photo: TRI LBM Team 2025, Figure 10, arXiv 2507.05331.}
  \note{[2 min] NARRATE: most evaluations compare a new checkpoint with the old one. Three requirements. Design: the two policies run interleaved and in random order on the same unit and fixture, and the operator who resets the cell and judges borderline runs does not know which checkpoint ran; otherwise winding heat and spring relaxation masquerade as a policy effect. Toyota Research Institute ran its large-behavior-model comparisons this way, as blind A/B trials. Test: with 4 failures among 40 runs, the chance that all 4 land in one arm is $\binom{20}{4}/\binom{40}{4} = 0.053$, two-sided 0.106. Neither checkpoint is shown better. Power: to separate 90 from 80 percent with 80 percent power takes about 195 runs per arm. ENGAGE: ``How many robot-learning papers report 20 trials per arm?'' Many; Snyder and colleagues show such gaps are often undetectable. FLEX: [OPTIONAL] comparison; land the segment here or on the previous slide: a clean record buys a bound, and each nine costs ten times the runs.}
\end{frame}

% =============================================================================
% SECTION: EXPOSURE WALL
% =============================================================================
\mlsysnote{[0.5 min] [CORE] Task success is counted per trial. Hazards are counted per hour. How many hours?}
\mlsysdivider{Can testing bound a rare hazard?}{Closed-Loop Evaluation, \S7.8 and \S7.9}
\section{Exposure Wall}

\begin{frame}{Bounding $10^{-9}$ hazards per hour needs 3.0 billion clean hours}
  \centering
  \mlsyssteps[width=0.95\textwidth,height=4.9cm,keepaspectratio]{exposure_step}{4}\\[0.1cm]
  {\normalsize $T \ge \ln(1/\asig)/\lambda \approx 3/\lambda$ is the rule of three in hours.}
  \mlsyscredit{Closed-Loop Evaluation, \S7.9.1, equation 9, Definition 7.2 and Napkin Math 7.3. One rig-year $=$ 8,766 h. Values from images/calcs.py.}
  \note{[2 min] NARRATE: a hazard is a rate per operating hour. With a constant rate $\lambda$ and no events in $T$ hours, the record's likelihood is $e^{-\lambda T}$; setting it to 0.05 gives the exposure wall, $T \ge 3/\lambda$. Step 1: the door's 99.9 percent trial claim, 33.3 hours. Step 2: a site requirement of $10^{-4}$ per hour needs 29,957 failure-free hours, 3.42 years of one machine. Step 3: $10^{-6}$ needs 3.0 million hours, 342 machine-years. Step 4: $10^{-9}$, aviation's catastrophic band, needs 3.0 billion hours: 100 rigs running continuously for 3,417 years (derived). The dashed line is one rig-year; it bounds the rate only to $3.4\times10^{-4}$ per hour, 342 times short of $10^{-6}$. WARN: parallel rigs add hours but not diversity; machines sharing a facility and a production batch bound the rate only for what they share, and simulated hours add no physical exposure. LINK: Lecture 15 applies the wall to injected faults. FLEX: [CORE] concept trace: exposure wall.}
\end{frame}

\begin{frame}{Do 56.7 million miles bound the fatality rate?}
  \mlsysactivity{Predict}{2 min}
  \begin{columns}[c]
    \begin{column}{0.30\textwidth}
      \centering\safeimg[width=\textwidth,height=3.4cm,keepaspectratio]{waymo.jpg}\\[0.05cm]
      {\footnotesize Waymo, San Francisco, 2023}
    \end{column}
    \begin{column}{0.68\textwidth}
      \centering
      {\normalsize\choice{A}{Yes}\quad\choice{B}{Only for injuries and crashes}}\\[0.25cm]
      \uncover<2->{\safeimg[width=\textwidth,height=4.2cm,keepaspectratio]{rand_waymo.pdf}}
    \end{column}
  \end{columns}
  \mlsyscredit{RAND RR-1478 (Kalra and Paddock 2016), Table 1, zero-failure rows; Kusano et al. 2025. Photo: Dllu, CC BY 4.0, Wikimedia Commons.}
  \note{[2 min] ENGAGE: vote. Step 2: B. RAND computed the miles a fleet must drive with zero events to bound its rate at the human benchmark with 95 percent confidence: 275 million miles for fatalities (1.09 per 100 million miles), 3.9 million for injuries, 1.6 million for crashes. Waymo's 56.7 million rider-only miles clear the injury and crash bars and fall short of the fatality bar by a factor of 4.9 (derived). Estimating the fatality rate within 20 percent needs 8.8 billion miles, 155 times Waymo's total. Kusano and colleagues could therefore measure what the exposure supports: 79 percent fewer injury crashes and 81 percent fewer airbag-deployment crashes than an exposure-matched human benchmark, with confidence intervals excluding zero. WARN: this is the opposite of the hook. Kusano's comparison matches road type, vehicle class, and city in the denominator; NHTSA's did not. Butler and Finelli reached the same wall for flight software in 1993. FLEX: [CORE].}
\end{frame}

\begin{frame}{Test arenas fix their conditions, and so bound their claims}
  \begin{columns}[c]
    \begin{column}{0.48\textwidth}
      \centering\safeimg[width=\textwidth,height=4.4cm,keepaspectratio]{nist-bays.jpg}\\[0.05cm]
      {\footnotesize NIST robot test bays, 2013}
    \end{column}
    \begin{column}{0.48\textwidth}
      \centering\safeimg[width=\textwidth,height=4.4cm,keepaspectratio]{mcity.jpg}\\[0.05cm]
      {\footnotesize Mcity test track, Michigan, 2023}
    \end{column}
  \end{columns}
  \mlsyscredit{NIST, Public domain, Wikimedia Commons. Mcity: USDA-FSA / NOAA orthophoto 2023, Public domain, Wikimedia Commons. Closed-Loop Evaluation, \S7.9.2, figure 6.}
  \note{[1.5 min] NARRATE: two real test facilities. Left: the NIST response-robot test bays, the same facility as the chapter's figure 6. Each station is a fixed obstacle with a standardized test method, so different robots can be compared run for run. Right: Mcity, a test track for automated vehicles, seen from the air: city blocks, intersections, and a highway loop built for repeatable trials. ENGAGE: ``What can a zero-failure run across every NIST bay not tell you?'' Expected: anything about glare, spills, pedestrians, or surfaces the arena does not contain. The control that makes tests repeatable also draws the edge of the evidence. A test campaign has to add those conditions on purpose, as declared strata. FLEX: [OPTIONAL] cut if short; the point is also made by the coverage grid.}
\end{frame}

\begin{frame}{A 12 N mean hides a 115 N peak past the latch limit}
  \centering
  \safeimg[width=0.95\textwidth,height=3.9cm,keepaspectratio]{mean_peak.pdf}\\[0.2cm]
  {\normalsize The first impact peak arrives faster than any reactive monitor.\\Cap the approach speed before contact at $v \le F_{\text{limit}}/\sqrt{m_{\text{red}}\,k}$.}
  \mlsyscredit{Closed-Loop Evaluation, \S7.9.2 and equation 10 (illustrative force values); ISO/TS 15066 Annex A.}
  \note{[2 min] NARRATE: average metrics mask structural overloads. Across one hundred door trials the mean latch contact force is 12 N, under the 15 N tripwire. When the gripper strikes the plate at a slight angle, the 99th-percentile peak reaches 115 N, past the 100 N structural limit. How can a 115 N peak occur under a 15 N tripwire? The first peak of an impact arrives a quarter oscillation after touch; if the monitor's sensing, compute, and actuation take longer than that rise time, no reactive software can attenuate it, however fast the model. Protection against the first peak comes before contact: passive compliance, which lowers stiffness and lengthens the rise, and a precontact speed cap such as the 0.03 m/s guarded approach. The bound is equation 10, the same relation behind ISO/TS 15066's transient contact limits. ENGAGE: ``Could twenty peak readings fit the tail beyond 100 N?'' Expected: no; extreme value theory, which Lecture 15 develops, needs many block maxima. FLEX: [OPTIONAL].}
\end{frame}

\begin{frame}{The evaluation record hands untested conditions to the runtime}
  \centering
  \mlsyssteps[width=0.98\textwidth,height=5.9cm,keepaspectratio]{record_step}{2}
  \mlsyscredit{Closed-Loop Evaluation, \S7.8 (table 5) and \S7.9.3 (five monitor channels); latch-bc-01 values from \S7.8.}
  \note{[3 min] NARRATE: this lecture's record, the link after the policy manifest in the record chain. Step 1: seven fields. The target ODD is the sampled cell, inside the manifest's declared ODD. Component hashes pin the manifest digest, the weights, and Unit 01's measured tolerances. The sampling protocol holds draw weights, seed, reset procedure, and the independence audit. The complete outcome log keeps every run on all channels. The statistical evidence states the method; twenty clean runs bound the failure rate at 13.9 percent with 95 percent confidence, and 299 would be needed for 99 percent. Step 2: the gap catalog lists what no trial visited: five cells, the loaded tote rack, other units. Each gap that can be sensed becomes a runtime monitor with a signal, a threshold, and a tested response, across five channels. A gap that cannot be sensed becomes a hard ODD limit or a physical barrier. LINK: Lecture 12's fallback ladder executes those responses; Lecture 16 reads this record into the safety case. FLEX: [CORE] concept trace: evaluation record.}
\end{frame}

\begin{frame}{Which monitor, if any, covers each gap?}
  \mlsysactivity{Diagnose}{2 min}
  \small
  \renewcommand{\arraystretch}{1.2}
  \begin{tabular}{@{}>{\raggedright\arraybackslash}m{0.45\textwidth}>{\raggedright\arraybackslash}m{0.50\textwidth}@{}}
    \toprule
    \textbf{Gap in latch-bc-01's record} & \uncover<2->{\textbf{Runtime obligation}} \\
    \midrule
    Strike plate displaced 3 mm & \uncover<2->{Input support (out-of-distribution score)} \\
    A chunk arrives after its 60 ms lease & \uncover<2->{Latency and freshness, then fallback} \\
    Dusty floor; wheels slip & \uncover<2->{Environmental assumption (wheel slip)} \\
    A clean glass partition in the aisle & \uncover<2->{\textcolor{RedLine}{\textbf{None can sense it:}} bumper, zoning, barrier} \\
    \bottomrule
  \end{tabular}
  \mlsyscredit{Closed-Loop Evaluation, \S7.9.3 (five channels; transduction failures).}
  \note{[2 min] ENGAGE: pairs fill the right column, ninety seconds. Step 2: a displaced plate is outside the training and test support, so the input-support monitor must flag it, with a measured false-negative rate. A late chunk is a latency and freshness event; the lease of Lecture 4 expires and a validated fallback takes over. Floor slip is an environmental assumption the slip ratio can check. The glass wall is the trap, because a transduction failure corrupts the telemetry itself, so every software monitor sees normal features and clear space. Only physical measures help: tactile bumpers, current limits, zoning, barriers, validated independently of the policy. WARN: a monitor that only logs provides no protection; each must command a response the permission path can execute. LINK: Lecture 18 calls this observational indistinguishability. Land the segment: past the exposure wall, safety evidence comes from the permission path and physical safeguards, and the record says where. FLEX: [CORE].}
\end{frame}

\begin{frame}{A humanoid fall needs its own denominator}
  \begin{columns}[c]
    \begin{column}{0.36\textwidth}
      \centering\safeimg[width=\textwidth,height=5.0cm,keepaspectratio]{atlas-testing.jpg}\\[0.05cm]
      {\footnotesize Atlas humanoid during testing, 2013.}
    \end{column}
    \begin{column}{0.60\textwidth}
      \small
      One failure can destroy the test article.\\[0.3cm]
      Count falls separately, with their own bound.\\[0.3cm]
      Declare disturbance strata first: spills, turf, curbs, pushes.\\[0.3cm]
      Human balance in demonstrations does not transfer to the robot.
    \end{column}
  \end{columns}
  \mlsyscredit{Photo: DARPA, Public domain, Wikimedia Commons. Closed-Loop Evaluation, Systems Perspective 7.2 (Evaluating a machine that falls).}
  \note{[1 min] NARRATE: the chapter's humanoid capstone. A failed guarded latch engagement leaves the wheeled machine intact for the next trial. When a biped loses balance, one failure can strike the floor and destroy actuators, links, and sensors. Photo: DARPA's Atlas humanoid during testing in 2013. Three rules follow. Falls are their own estimand, with their own denominator and bound, never folded into task completion. Disturbance strata (low-friction spills, turf, curbs, pushes) and their weights are declared before testing; polished concrete supports nothing else. Retargeted human motion carries the human's balance instead of the robot's, so missing recovery states go in the gap catalog and are investigated under containment. FLEX: [OPTIONAL] capstone; skip if short.}
\end{frame}

% =============================================================================
% SECTION: WRAP-UP
% =============================================================================
\mlsysnote{[0.5 min] [CORE] Check the ideas, then return to the 40 percent.}
\mlsysdivider{What does a clean record support?}{Closed-Loop Evaluation, \S7.11 and \S7.12}
\section{Wrap-Up}

\begin{frame}{Vote: true or false?}
  \mlsysactivity{Vote}{2 min}
  \small
  \renewcommand{\arraystretch}{1.35}
  \begin{tabular}{@{}>{\raggedright\arraybackslash}m{0.48\textwidth}>{\raggedright\arraybackslash}m{0.46\textwidth}@{}}
    \toprule
    \textbf{Claim} & \uncover<2->{\textbf{Verdict}} \\
    \midrule
    Zero failures in 1,000 trials proves 99.9 percent. & \uncover<2->{\textbf{\textcolor{RedLine}{False.}} The lower limit is 99.70\%; 2,995 needed.} \\
    Zero collisions under supervision proves the policy reliable. & \uncover<2->{\textbf{\textcolor{RedLine}{False.}} Interventions censored the crashes.} \\
    Lower offline action error means better rollouts. & \uncover<2->{\textbf{\textcolor{RedLine}{False.}} The lowest-loss checkpoint scored 2.7\%.} \\
    Stopping once the tally looks good keeps 95 percent confidence. & \uncover<2->{\textbf{\textcolor{RedLine}{False.}} Only a boundary set in advance.} \\
    \bottomrule
  \end{tabular}
  \mlsyscredit{Closed-Loop Evaluation, \S7.11 (fallacies and pitfalls); Mandlekar et al. 2021, Appendix G.}
  \note{[2 min] ENGAGE: students write T or F for each, then compare with a neighbor. Step 2: all four are false, each a fallacy or pitfall from the chapter. One: $0.05^{1/1000} = 0.9970$, short of 99.9 percent; 2,995 clean runs meet it (derived). Two: zero damage under active supervision records a sequence of interventions; report ACR and MTBI beside it. Three: offline error measures imitation on logged states; the lowest-loss checkpoint succeeded 2.7 percent of the time against 80.7 for the highest-scoring one. Four: peeking inflates false acceptance; a sequential test declared before the first run controls it. If many students marked the first true, return to the focus slide. FLEX: [CORE].}
\end{frame}

\begin{frame}{The denominator decided the 40 percent claim}
  \begin{columns}[c]
    \begin{column}{0.30\textwidth}
      \centering\safeimg[width=\textwidth,height=3.6cm,keepaspectratio]{nhtsa_fig11.png}
    \end{column}
    \begin{column}{0.67\textwidth}
      \small
      \renewcommand{\arraystretch}{1.4}
      \begin{tabular}{@{}>{\raggedright\arraybackslash}p{0.29\textwidth}>{\raggedright\arraybackslash}p{0.27\textwidth}>{\raggedright\arraybackslash}p{0.34\textwidth}@{}}
        \toprule
        \textbf{Question} & \textbf{NHTSA 2017} & \textbf{Door campaign} \\
        \midrule
        Which exposure? & All miles & 20 runs, one cell \\
        Who acted? & Not recorded & All interventions kept \\
        What bound? & None stated & Failure $\le$ 13.9\% \\
        What is left? & Unstated & Five cells, payload, units \\
        \bottomrule
      \end{tabular}
    \end{column}
  \end{columns}
  \mlsyscredit{NHTSA PE 16-007, Figure 11. Closed-Loop Evaluation, War Story 7.1 and \S7.8.}
  \note{[2 min] NARRATE: return to the hook. The 40 percent drop divided airbag crashes by all miles, engaged or not, for vehicles whose installation mileage was often missing; it recorded no engagement, no ODD, and stated no bound. The door campaign answers the same four questions. Its exposure is twenty runs in one cell of six. Every clip and takeover stays in the log, credited to whoever acted. Its bound is exact, a failure rate of at most 13.9 percent at 95 percent confidence, for Unit 01, nominal plate, guarded approach, empty rack. What it leaves is named: five cells, the loaded machine, other units, each handed to a monitor or a hard limit. ENGAGE: ``What would NHTSA have needed to report?'' Expected: miles with Autosteer engaged, by road type, with complete exposure records and a confidence interval. FLEX: [CORE] this slide answers the hook.}
\end{frame}

\begin{frame}{Takeaways}
  \thumbline{thumb_coverage_step2.pdf}{A success rate describes a declared population of trials on one machine; interventions count against the policy.}\\[0.15cm]
  \thumbline{thumb_aisle_step2.pdf}{Offline scores reset each step to recorded states, so they miss drift; each evidence regime has its own blind spot.}\\[0.15cm]
  \thumbline{thumb_inversion_step2.pdf}{Twenty clean runs bound failure at 13.9 percent; each added nine costs ten times the trials.}\\[0.15cm]
  \thumbline{thumb_record_step2.pdf}{Rare hazards need about $3/\lambda$ clean hours, so the evaluation record hands every untested condition to the runtime.}
  \note{[1 min] [CORE] One line per segment. Ask students which slide taught each line, and to star the one they would have gotten wrong before today. Name the record: the evaluation record completes the first strand of the record chain, pairing the plant limits of the limit record with the bounded competence recorded here.}
\end{frame}

\begin{frame}{Open questions}
  \begin{columns}[c]
    \begin{column}{0.62\textwidth}
      \small
      \begin{itemize}\setlength{\itemsep}{0.2cm}
        \item A fleet logs millions of hours, all in one facility. How much diversity must a claim sample, and who decides?
        \item Learned world models now generate test scenes. What paired real evidence should earn one the right to screen?
        \item Interventions censor the failures a test most needs. Should some trials run to failure on sacrificial hardware?
      \end{itemize}
    \end{column}
    \begin{column}{0.35\textwidth}
      \centering\safeimg[width=\textwidth,height=3.8cm,keepaspectratio]{mcity.jpg}\\[0.05cm]
      {\footnotesize Mcity, aerial, 2023}
    \end{column}
  \end{columns}
  \mlsyscredit{Mcity: USDA-FSA / NOAA orthophoto 2023, Public domain, Wikimedia Commons.}
  \note{[1 min] [OPTIONAL] No single answer. Use one for a two-minute discussion or a short written assignment. The first comes from the exposure wall's pitfall, that parallel rigs add hours but no diversity. The second extends the simulator rule of \S7.4.2 to learned world models. The third weighs irreversibility against the censoring problem; the humanoid capstone is one side of it.}
\end{frame}

\begin{frame}{Read and explore}
  \small
  \begin{itemize}\setlength{\itemsep}{0.04cm}
    \item \textbf{Read:} Closed-Loop Evaluation, \S7.1 to \S7.12; Checkpoints 7.1 to 7.3; redo Napkin Math 7.2 for 99.5 percent.
    \item Hanley and Lippman-Hand, ``If Nothing Goes Wrong, Is Everything All Right?'', JAMA 1983.
    \item Kalra and Paddock, \emph{Driving to Safety}, RAND RR-1478, 2016.
    \item Butler and Finelli, ``The Infeasibility of Quantifying the Reliability of Life-Critical Real-Time Software,'' IEEE TSE 1993.
    \item Kress-Gazit et al., ``Robot Learning as an Empirical Science: Best Practices for Policy Evaluation,'' 2024.
    \item Li et al., ``Evaluating Real-World Robot Manipulation Policies in Simulation'' (SIMPLER), CoRL 2024.
  \end{itemize}
  \note{[0.5 min] [CORE] The chapter first, with its three checkpoints. Each outside reading backs one idea: the rule of three (Hanley), the exposure wall in road miles (Kalra and Paddock), the same wall for flight software (Butler and Finelli), the trial protocol as a reporting standard (Kress-Gazit), and paired sim-and-real screening (Li). FLEX: add the evaluation lab your course uses.}
\end{frame}

\begin{frame}[c]
  \begin{columns}[c]
    \begin{column}{0.34\textwidth}
      \centering\mlsyskitfig[width=\textwidth,height=4.2cm,keepaspectratio]{stack_boundary_brain}
    \end{column}
    \begin{column}{0.62\textwidth}
      {\small\textcolor{lightgray}{NEXT LECTURE \textbar{} SENSOR PERCEPTION}}\\[0.5cm]
      {\Large\bfseries A detection can be correct and already late.\\[0.15cm] What must it state?}
    \end{column}
  \end{columns}
  \note{[0.5 min] [CORE] Part II ends with a record that confines trust in the learned proposer to its target ODD and names the monitors the runtime must supply. Lecture 8 opens Part III and the third law, proposal is not permission. Its first obligation falls on perception, because an input-support or latency monitor can act only on observations that carry their capture time, their frame, and their measured error, while the base keeps moving during exposure, readout, and inference.}
\end{frame}

\end{document}
